Ring homomorphisms send nilpotent elements to nilpotent elements, so the nilradicals form a direct system. By \exref{2.19}, their inclusions into $A_i$ induce an injection $\varinjlim_i\operatorname{nil}(A_i)\to A$, where $A=\varinjlim_iA_i$. Its image consists of nilpotent elements.

Conversely, if $\alpha_i(x)^n=0$, then $\alpha_i(x^n)=0$. By \exref{2.15}, some $j\geq i$ satisfies $\alpha_{ij}(x)^n=0$. Thus $\alpha_i(x)=\alpha_j(\alpha_{ij}(x))$ comes from the nilradical of $A_j$.

If each $A_i$ is an integral domain, then $A\neq0$ by \exref{2.21}. Represent two elements at a common stage as $\alpha_i(x),\alpha_i(y)$. If their product is zero, then $\alpha_{ij}(x)\alpha_{ij}(y)=0$ at some later stage by \exref{2.15}. Since $A_j$ is a domain, one factor is zero, so one of the original elements is zero.
